\section{Methodology}

This section outlines the comprehensive methodology employed in analyzing the ESG performance of Indonesian companies. The approach combines quantitative data analysis, qualitative assessment, and contextual understanding of Indonesia's unique ESG landscape.

\subsection{Data Collection and Sources}

The primary data sources for this analysis include:

1. Annual reports and sustainability reports of Indonesian companies
2. ESG ratings from major providers such as Sustainalytics, MSCI, and ISS ESG
3. Government databases and reports from Indonesian regulatory bodies
4. Industry-specific reports and academic publications

Data was collected for the fiscal year 2022, with a focus on companies listed on the Indonesia Stock Exchange (IDX) across various sectors.

\subsection{ESG Framework and Indicators}

The analysis follows a comprehensive ESG framework tailored to the Indonesian context. The framework is based on the Global Reporting Initiative (GRI) Standards \cite{GRI2021} and incorporates elements from the Sustainability Accounting Standards Board (SASB) \cite{SASB2022} and the Task Force on Climate-related Financial Disclosures (TCFD) \cite{TCFD2021}.

The framework is structured as follows:

\begin{itemize}
    \item \textbf{Environmental (E)}:
    \begin{itemize}
        \item Carbon emissions and energy consumption
        \item Water usage and management
        \item Waste management and circular economy initiatives
        \item Biodiversity conservation efforts
        \item Climate change mitigation and adaptation strategies
    \end{itemize}
    \item \textbf{Social (S)}:
    \begin{itemize}
        \item Labor practices and human rights
        \item Community engagement and development
        \item Health and safety measures
        \item Diversity and inclusion initiatives
        \item Supply chain management
    \end{itemize}
    \item \textbf{Governance (G)}:
    \begin{itemize}
        \item Board structure and diversity
        \item Executive compensation
        \item Anti-corruption measures
        \item Shareholder rights
        \item Risk management and internal controls
    \end{itemize}
\end{itemize}

\subsection{Quantitative Analysis}

For each ESG category, key performance indicators (KPIs) were identified and quantified where possible. These KPIs were then normalized and weighted based on their relative importance in the Indonesian context.

For example, in the Environmental category, carbon emissions intensity was calculated as CO2 equivalent emissions per unit of revenue. This was then compared to industry benchmarks and national targets set by the Indonesian government \cite{IndonesiaNDC2021}.

In the Social category, employee turnover rates and workplace accident frequencies were analyzed to assess labor practices. These metrics were benchmarked against industry averages and best practices in Southeast Asia.

For Governance, board independence ratios and the presence of women on boards were examined. The analysis also considered the implementation of the Indonesian Corporate Governance Code \cite{IDXCGC2020}.

\subsection{Qualitative Assessment}

In addition to quantitative analysis, a qualitative assessment was conducted to evaluate aspects of ESG performance that are not easily quantifiable. This included:

1. Review of corporate policies and commitments related to ESG issues
2. Analysis of stakeholder engagement practices
3. Assessment of transparency and disclosure quality
4. Evaluation of ESG integration into business strategy

The qualitative assessment was based on a scoring system that considered the comprehensiveness, specificity, and ambition of ESG-related disclosures and initiatives.

\subsection{Contextual Analysis}

Given Indonesia's unique ESG challenges and opportunities, the analysis incorporated a contextual layer. This included:

1. Assessment of companies' alignment with Indonesia's National Medium-Term Development Plan (RPJMN) \cite{Bappenas2020}
2. Evaluation of companies' contributions to the United Nations Sustainable Development Goals (SDGs) \cite{UNSDGS2022}
3. Consideration of Indonesia's commitments under the Paris Agreement and its Nationally Determined Contributions (NDCs) \cite{IndonesiaNDC2021}
4. Analysis of companies' responses to Indonesia-specific ESG issues such as deforestation, peatland management, and marine ecosystem conservation

\subsection{Weighting and Aggregation}

The final ESG scores were calculated by aggregating the weighted scores across all indicators. The weighting scheme was developed based on:

1. The relative importance of each ESG category to the company's sector
2. The materiality of ESG issues in the Indonesian context
3. Expert consultation with ESG specialists familiar with the Indonesian market

The aggregation formula used is as follows:

\begin{equation}
ESG\_Score = w_E \times E\_Score + w_S \times S\_Score + w_G \times G\_Score
\end{equation}

where $w_E$, $w_S$, and $w_G$ are the weights for Environmental, Social, and Governance categories respectively, and $E\_Score$, $S\_Score$, and $G\_Score$ are the normalized scores for each category.

\subsection{Limitations and Future Improvements}

This methodology, while comprehensive, has certain limitations:

1. The reliance on publicly available data may result in incomplete information for some companies.
2. The qualitative assessment introduces a degree of subjectivity.
3. The weighting scheme, while based on expert consultation, may not fully capture the nuances of every industry.

Future improvements to the methodology could include:

1. Incorporating more granular, company-specific data through direct engagement
2. Developing industry-specific ESG frameworks to better capture sector-specific risks and opportunities
3. Integrating emerging ESG issues such as biodiversity loss and social inequality more explicitly into the analysis

\subsection{Conclusion}

This methodology provides a robust framework for assessing the ESG performance of Indonesian companies. By combining quantitative analysis, qualitative assessment, and contextual understanding, it offers a comprehensive view of companies' ESG practices and their alignment with national and global sustainability goals. The approach is designed to be adaptable, allowing for continuous refinement as ESG considerations evolve in the Indonesian context.